\chapter{Discussion and Limitations}
\label{ch:discussion}

\section{What Worked}

Binary detection on the day split survives temporal shift better than
prior expectations. Random forest at F1 = 0.86 with a 6.5\% FPR is the
deployment-relevant ballpark for organisations that can absorb a
moderate alert budget --- though production deployment requires
live-traffic validation, drift monitoring, and the network-layer
mitigations discussed in Chapter~\ref{ch:adversarial}. The LODO study
gave a stronger generalisation claim than any single split: ROC-AUC =
$0.926 \pm 0.053$ across days with very different attack mixes is a
narrow band that survives genuinely disjoint test scenarios. The MITRE
ATT\&CK mapping turns probability scores into operational alerts. The
reproducible pipeline (45--60 minutes, single Makefile target) is a
real artefact a successor researcher can extend.

\paragraph{Engineering maturity.}
Three small engineering choices accumulated into the most useful
property of this work: a deterministic \texttt{random\_state}, a single
Makefile entry point that runs the whole pipeline end to end, and an
explicit leakage-check module that fails loudly if the four rate
features are reintroduced. None of these are research contributions ---
they are the difference between a thesis a successor researcher can
extend in an afternoon and one they have to rebuild from scratch.

\section{What Did Not Work}

Multi-class classification on the day split. Naïve adversarial training.
Both negative results are informative: the first establishes a hard
structural limit (the model cannot output a label that does not exist in
training), the second establishes that off-the-shelf defences fail. The
adversarial-training failure is consistent with Tramer et
al.~\cite{Tramer2020}, who documented similar failures across many
adversarial-training schemes. Reporting these failures rather than
hiding them is part of the contribution: a successor researcher knows
what not to waste time on. The naïve defence dropped clean F1 from
0.859 to 0.734 and increased evasion from 24.4\% to 26.8\% --- the
worst possible outcome for a defence.

Three findings during the experimental phase surprised the author and
are worth recording. The most concentrated adversarial brittleness:
83.6\% of successful evasions move a single feature, Flow IAT Min. The
narrow LODO ROC-AUC variance for random forest at $\pm 0.053$ across
genuinely disjoint test scenarios. And the sharp rank reversal between
CICIDS and UNSW datasets: the CICIDS-best model is not the UNSW-best.
Headline benchmark numbers measure the dataset they were generated on at
least as much as they measure the quality of the model.

\section{Threats to Validity}

\paragraph{Dataset limitations.} Both CICIDS2017 and UNSW-NB15 are
synthetic captures from controlled testbeds. Real enterprise traffic is
messier: more diverse user behaviour, more encrypted (TLS) traffic, a
long-tailed distribution of rare attack variants. Numbers obtained on
synthetic data are likely upper bounds on what would be achievable in
production. CICIDS2017 also has documented labelling
errors~\cite{Engelen2021} which we did not relabel.

\paragraph{Evaluation limitations.} LODO is over a single dataset;
cross-dataset LODO would be stronger. Multi-threaded non-determinism
introduces approximately 1\% F1 drift between runs. McNemar's test
assumes independent samples; flows in CICIDS share temporal
correlations, so the McNemar $p$-values reported in
Section~\ref{sec:mcnemar} should be read as supporting evidence for the
model ranking rather than as exact significance levels.

\paragraph{Attacker-model limitations.} The threat model assumes
score-query black-box access with unlimited queries; many real
deployments rate-limit query access, which slows but does not change the
qualitative outcome. The $L^\infty$ budget on a 19-feature subset is a
simplification. Most importantly, the thesis perturbs feature vectors
rather than synthesising adversarial network packets. Bridging the
problem-space gap (Pierazzi et al.~\cite{Pierazzi2020}) is open work.

\paragraph{Deployment limitations.} The MITRE mapping is a static JSON
file rather than a learned mapping over richer telemetry. There is no
streaming or online mode. There is no adversary simulator: we assume
the attacker is greedy and uses a fixed budget, while a sophisticated
real attacker would adapt their strategy based on observed detector
behaviour.

\section{Relation to the Published Literature}

The thesis sits in a small but growing literature on honest evaluation
of ML-based intrusion detection. Engelen et al.~\cite{Engelen2021}
flagged the leakage rate features in CICIDS2017 --- we replicate the
finding and drop those columns explicitly, confirming the
four-percentage-point effect on stratified RF macro F1. Sharafaldin et
al.~\cite{Sharafaldin2018} reported macro F1 $\approx 0.97$ on CICIDS
using stratified splits; we reproduce 0.86 (XGBoost) on the same
protocol after leakage filtering, and 0.44 on the day split. Apruzzese
et al.~\cite{Apruzzese2019} demonstrated adversarial brittleness for
flow-based IDS classifiers; the 21.4\% bounded evasion rate
(24.4\% unbounded headline) reported here is in the same ballpark, with
the additional finding that the brittleness is concentrated in a single
feature. The pattern across the references is consistent: the attack
side of the literature is mature, and the defence side has many more
unsolved problems than solved ones.
